\section*{Bab \ref{ch:b2:vertebrate-development} --- Perkembangan Lembaga Vertebrata}

\begin{solution}{exo:b2:vertebrate-development:1}
Pembelahan: mitosis cepat tanpa pertumbuhan, yang menghasilkan
\omterm{def:b2:vertebrate-development:blastula}{blastula}. \omterm{def:b2:vertebrate-development:blastula}{Blastula}: sebuah bola atau cakram berongga berisi sel kecil
mengelilingi \omterm{def:b2:vertebrate-development:blastula}{blastosolnya}. Gastrulasi: gerakan sel yang membawa
mesoderm dan endodermnya masuk, yang menghasilkan gastrula berlapis
tiga dengan sebuah rongga usus dan sebuah notokord. Neurulasi:
pelipatan ektoderm punggungnya menjadi \omterm{prop:b2:vertebrate-development:neurulation}{tabung saraf}, beserta
peruasan somitnya, yang menghasilkan neurula dengan rancang bangun
tubuh vertebrata.
\end{solution}

\begin{solution}{exo:b2:vertebrate-development:2}
Ektoderm: epidermis, otak dan sumsum tulang belakang, saraf tepi (\omterm{prop:b2:vertebrate-development:neurulation}{pial
saraf}), lensa, dan sel pigmen. Mesoderm: notokord, otot, tulang dan
rawan, ginjal, jantung dan darah, dermis. Endoderm: pelapis ususnya,
hati, pankreas, paru-paru, kelenjar tiroid.
\end{solution}

\begin{solution}{exo:b2:vertebrate-development:3}
Sedikit \omterm{def:b2:vertebrate-development:egg}{kuning telur}: pembelahan lengkap yang nyaris sama besar dan
gastrulasi lewat involusi menembus sebuah blastopor bulat di dalam
sebuah bola berongga (katak: lengkap tetapi tidak sama besar, dengan
sel tumbuhan berkuning telur yang besar dan lambat). Banyak \omterm{def:b2:vertebrate-development:egg}{kuning
telur} (anak ayam): pembelahannya terbatas pada sebuah cakram di atas
\omterm{def:b2:vertebrate-development:egg}{kuning telurnya}, dan gastrulasinya lewat sebuah celah, yaitu garis
primitif, di dalam cakram datarnya, dengan lapisannya menyebar di atas
\omterm{def:b2:vertebrate-development:egg}{kuning telurnya} alih-alih mengurung sebuah rongga.
\end{solution}

\begin{solution}{exo:b2:vertebrate-development:4}
\omterm{def:b2:vertebrate-development:induction}{Pengimbasan} adalah pergantian nasib sebuah jaringan yang berkompeten
oleh sebuah isyarat dari jaringan tetangganya. Sel tumbuhannya mengimbas
mesoderm di dalam zona tepi di atasnya; \omterm{def:b2:vertebrate-development:induction}{pengaturnya} (mesoderm bibir
punggung) mengimbas lempeng saraf di dalam ektoderm punggungnya;
gelembung optiknya mengimbas lensa di dalam ektoderm kepalanya.
\end{solution}

\begin{solution}{exo:b2:vertebrate-development:5}
$2^{12} = 4096$ sel berjari-jari $1/2^{4} = \qty{62.5}{\micro m}$;
permukaannya dikalikan $2^{4} = 16$. Permukaan tambahannya adalah
membran epitelium yang akan dilipat gastrulasinya menjadi kulit dan
usus, serta luas yang lewatnya \omterm{def:b2:vertebrate-development:blastula}{blastulanya} bertukar dengan
sekelilingnya.
\end{solution}

\begin{solution}{exo:b2:vertebrate-development:6}
$10^{-4}\times 2^{k} = 0.4$: $2^{k} = 4000$, $k = 12$. Dengan DNA empat
kali lipat, rasionya bermula pada $4\times 10^{-4}$ lalu mencapai $0.4$
pada $2^{k} = 1000$: yaitu pembelahan kesepuluh, dua lebih awal.
\end{solution}

\begin{solution}{exo:b2:vertebrate-development:7}
$k = \ln 2/600 = \qty{1.16e-3}{s^{-1}}$; $\lambda = \sqrt{2\times
10^{-12}/1.16\times 10^{-3}} = \qty{42}{\micro m}$. Ambangnya pada
$x = \lambda\ln 2 = \qty{29}{\micro m}$ dan $\lambda\ln 10 =
\qty{96}{\micro m}$: pita selebar 29, 67, dan sisanya.
\end{solution}

\begin{solution}{exo:b2:vertebrate-development:8}
Melipatduakan sumbernya menggeser kedua batasnya ke luar sejauh
$\lambda\ln
2 = \qty{29}{\micro m}$: pita pertamanya berlipat dua, pita
keduanya mempertahankan lebarnya. Melipatduakan $k$ membagi $\lambda$
dengan $\sqrt 2$ (\qty{29}{\micro m}) dan, karena $c_0 = j/\sqrt{Dk}$,
membagi $c_0$ dengan $\sqrt 2$ pula: batasnya bergerak ke
$\lambda'\ln(c_0'/T) =
29\times(0.69 - 0.35) = \qty{10}{\micro m}$ dan
$29\times(2.30 -
0.35) = \qty{57}{\micro m}$ --- kedua pitanya menyusut.
\end{solution}

\begin{solution}{exo:b2:vertebrate-development:9}
Sebuah bibir punggung dari gastrula kadal air tanpa pigmen dicangkokkan
ke sisi perut gastrula yang berpigmen; ia bergulir masuk dan inangnya
membentuk sebuah lembaga kedua --- \omterm{prop:b2:vertebrate-development:neurulation}{tabung saraf}, somit, usus --- yang
bersambung dengan yang pertama. Perbedaan pigmennya menunjukkan bahwa
sumbu kedua itu terbuat dari sel inangnya, sehingga cangkokannya telah
mengimbas mereka alih-alih membangun sumbunya sendiri. Hal itu
menyingkirkan pendiferensiasian sendiri cangkokannya sebagai
penjelasannya, dan menunjukkan bahwa sebuah daerah kecil mengusung
isyarat yang mengatur seluruh sumbunya.
\end{solution}

\begin{solution}{exo:b2:vertebrate-development:10}
(a) Kulit perut: ektoderm dininya berkompeten tetapi belum terimbas,
sehingga ia mengikuti sekeliling barunya. (b) Jaringan saraf: pada
gastrula lanjut ia telah terimbas dan terdeterminasi. (c) Lempeng
saraf: ektoderm perut gastrula lanjut masih berkompeten, dan
notokordnya mengimbasnya. (d) Tanpa struktur punggung: sebuah bola
jaringan perut, karena tidak ada yang mengimbas mesodermnya menuju
nasib punggung maupun ektodermnya menuju nasib saraf.
\end{solution}

\begin{solution}{exo:b2:vertebrate-development:11}
Pada \qty{18}{\celsius}: $90 + 11\times 30 = \qty{420}{min}$, yaitu 7
jam. Pada \qty{10}{\celsius}: \qty{17.5}{h}. Peralihannya jatuh pada
pembelahan kedua belas pada kedua kasusnya karena yang memicunya adalah
rasio DNA terhadap sitoplasmanya, yang bergantung pada cacah
pembelahannya dan bukan pada berapa lama pembelahannya berlangsung:
jamnya mencacah pelipatduaan, bukan jam.
\end{solution}

\begin{solution}{exo:b2:vertebrate-development:12}
Pada perkembangan regulatif, selnya memperoleh nasibnya lewat
pengimbasan dari tetangganya, sehingga lembaganya dapat membangun
kembali sesudah kehilangan (pengembaran dari telur yang terbelah,
regenerasi daerah yang dicangkokkan); pada perkembangan mosaik,
nasibnya diwarisi bersama penentu sitoplasma yang terlokalkan, sehingga
tiap blastomernya hanya membuat bagiannya dan lembaga yang terbelah
membuat dua setengah larva. Perbedaannya hanyalah derajat: sabit kelabu
kataknya adalah sebuah penentu, dan lembaga mosaik memakai pengimbasan
kemudian; tiap lembaganya mencampur pewarisan dan percakapan, dan
vertebrata bersandar pada percakapannya lebih lama.
\end{solution}

\begin{solution}{pb:b2:vertebrate-development:1}
\textbf{1.} Telurnya $\tfrac43\pi(0.7)^{3} = \qty{1.44}{mm^3}$;
nukleusnya $\tfrac43\pi(0.03)^{3} = \qty{1.1e-4}{mm^3}$; rasionya
$8\times
10^{-5}$.
\textbf{2.} $2^{k} = 4000$: $k = 12$, yaitu 4096 sel.
\textbf{3.} $75 + 11\times 30 = \qty{405}{min}$, kira-kira 6,75 jam.
\textbf{4.} Jari-jarinya $0.7/2^{4} = \qty{44}{\micro m}$; rasionya
$8\times
10^{-5}\times 4096 = 0.32$ --- yaitu rasio sebuah sel biasa:
nukleusnya telah menyusul sitoplasmanya.
\textbf{5.} $2^{4} = 16$.
\textbf{6.} Kira-kira $4095\,c$ DNA; telurnya memuat cukup histon,
polimerase, dan nukleotida tersimpan bagi dua belas pembelahan,
sehingga tidak ada transkripsi yang dibutuhkan.
\textbf{7.} Gen zigotiknya menyala, pembelahannya melambat dan
kehilangan keserentakannya, dan selnya menjadi mampu bergerak.
Menyuntikkan DNA berlebih memajukan peralihannya sebanyak pembelahan
seperti berlebihnya DNA itu, yang tidak dapat dijelaskan sebuah
pencacah pembelahan maupun pencacah waktu.
\textbf{8.} Berseberangan dengan titik masuk spermanya, pada sabit
kelabu yang dipancangkan perputaran korteks telurnya sesudah
pembuahannya.
\textbf{9.} Mesoderm punggungnya lebih dahulu (lempeng prakorda, lalu
notokord) ke atap arkenteronnya pada garis tengahnya; mesoderm
sampingnya (somit, lempeng samping) di sebelahnya; endodermnya melapisi
arkenteronnya; ektodermnya menyebar menutupi bagian luarnya lewat
epiboli.
\textbf{10.} $\qty{2000}{\micro m}/\qty{600}{min} = \qty{3.3}{\micro
m/min}$: tiga kali sebuah fibroblas --- sebuah lembaran yang terpadu
bergerak lebih cepat daripada sel yang sendirian.
\textbf{11.} Gerakan selnya membutuhkan hasil gen zigotiknya (molekul
pelekat baru, kemampuan bergerak) dan daur yang melambat; sebelum
peralihannya, selnya hanya dapat membelah.
\textbf{12.} Tiap kordata memiliki sebuah notokord pada suatu tahapnya,
dan vertebrata membangun tiang tulang belakangnya di sekelilingnya;
\omterm{prop:b2:vertebrate-development:neurulation}{tabung sarafnya} diimbas olehnya dan merupakan akibat, bukan asal,
rancangannya.
\textbf{13.} Tanpa pengerutan sel botolnya tidak ada blastopor yang
terbentuk, tidak ada yang bergulir masuk, dan lembaganya tetap sebuah
bola dengan lapisannya di tempatnya: tanpa usus, tanpa notokord, tanpa
sumbu.
\textbf{14.} $k = \ln 2/1200 = \qty{5.8e-4}{s^{-1}}$; $\lambda =
\sqrt{10^{-12}/5.8\times 10^{-4}} = \qty{42}{\micro m}$.
\textbf{15.} $x_1 = 42\ln 2.5 = \qty{38}{\micro m}$; $x_2 = 42\ln 12.5
= \qty{105}{\micro m}$.
\textbf{16.} Pita selebar \qty{38}{\micro m} (2,5 sel), \qty{67}{\micro
m} (4,5 sel), dan sisanya.
\textbf{17.} $x_1 = 42\ln 1.25 = \qty{9}{\micro m}$; $x_2 = 42\ln 6.25
= \qty{77}{\micro m}$: kedua batasnya bergerak ke dalam sejauh
$\lambda\ln 2 =
\qty{29}{\micro m}$; pita pertamanya menyusut dari 38
menjadi 9, pita tengahnya mempertahankan lebarnya, dan pita luarnya
tumbuh.
\textbf{18.} Kalau dipindahkan sebelum pembacaannya, ia menjadi nasib A
--- ia membaca kepekatan di tempatnya yang sekarang; kalau dipindahkan
sesudahnya, ia mempertahankan nasib B --- ia sudah terdeterminasi.
\textbf{19.} $x_T = \lambda\ln(c_0/T)$: perubahan sumbernya dengan
faktor $f$ menggeser tiap batasnya sejauh $\lambda\ln f$, sehingga
galat dua kali lipat hanya menggeser polanya sejauh $0.7\lambda$, yaitu
sebagian kecil sebuah medan yang merentang beberapa $\lambda$ ---
polanya nyaris tak peka terhadap jumlah morfogennya yang tepat.
\textbf{20.} Bibir punggung ke perut: sebuah sumbu kedua. Zona tepi
perut ke punggung: ia ditetapkan ulang oleh sekelilingnya lalu membuat
jaringan punggung; tidak ada tambahan apa pun.
\textbf{21.} Pada sebuah neurula, ektodermnya tidak lagi berkompeten:
cangkokannya membuat notokord tetapi tidak mengimbas \omterm{prop:b2:vertebrate-development:neurulation}{tabung saraf} kedua.
\textbf{22.} Dibelah sepanjang bidang sabitnya: dua lembaga yang
lengkap; dibelah tegak lurus: satu lembaga dan sebuah bola jaringan
perut.
\textbf{23.} Sendirian: sebuah bola epidermis. Bersama sel tumbuhannya:
mesoderm --- otot, notokord --- yang terimbas di dalam tudungnya.
\textbf{24.} Untuk membuktikan bahwa sumbu keduanya diimbas di dalam
sel inangnya dan bukan dibangun cangkokannya; kalau cangkokannya tak
terbedakan, pendiferensiasian sendiri tidak dapat disingkirkan. Sebuah
cangkokan dari spesies yang berbeda berhasil (mereka memakai dua
spesies kadal air), dan itulah yang membuat selnya dapat dibedakan.
\textbf{25.} Pembelahan kedua belas, 4096 sel, pada kira-kira 6,75 jam;
pita selebar 38 dan \qty{67}{\micro m}, serta segala sesuatu di luar
\qty{105}{\micro m}.
\end{solution}
